\scanpage{36}

Soit

\[
\unsure{\mathcal A^{*}|\mathcal A^{\times}|A}=\deleted{\operatorname{Aff}}\operatorname{Aut}(T,K(\tfrac12))\deleted{=\operatorname{Aff}\unsure{(}}
\tag*{(38)}
\]

donc

\[
1\longrightarrow\deleted{\unsure{\mu_\infty}}T\mathrel{\unsure{-|\longrightarrow}}\mathcal A\longrightarrow
\underset{\substack{\wr\\\operatorname{Aut}(T)}}{\widehat{\mathbb Z}^{\times}}\mathrel{-}1.
\tag*{(39)}
\]

\deleted{Les \inserted{groupes $\mu_n,\mu_\infty=T$,} ensembles $\Pi_n,\Pi_\infty$, et les ensembles $\underline E(I),E(I)$ \unsure{soit} \unsure{d'où un hom} \inserted{$\underline E(\alpha_*),E(\alpha_*)$}}

\[
\deleted{\Gamma\longrightarrow\mathcal A}
\]

$\mathcal A$ opère sur $\Pi_n,\Pi_\infty\inserted{=T^3/T}$ par l'intermédiaire de $\widehat{\mathbb Z}^{*}$, (par homothéties \unsure{adiques|adéquates} \unsure{..|...}) et le groupe

\[
\widetilde{\mathcal E}=(\unsure{\mathfrak S_3|\mathscr S_3}\times\mathcal A)\underset{1/2}{\cdot}\Pi_\infty
\]

% slanted
\margin{$\alpha(0)=0$\\ $\alpha(1)=1$\\ $\alpha(\infty)=\infty$\\ $\alpha(*)=\{\alpha(0),\alpha(1),\alpha(\infty)\}$}

opère \inserted{de façon évidente} sur les ensembles $\underline E_n(\alpha\unsure{(*)|{}_*}),E(\alpha_*))$ $\underline E(\unsure{\sigma|\alpha|\varpi}_*),E(\sigma_*),\underline E(\deleted{\beta_*}\alpha'_*),E(\deleted{\beta_*}\alpha'_*)$, -- en fait sur les quatre premiers il opère via le quotient $\deleted{\widetilde{\mathcal E}/T\simeq}\unsure{\widehat{\mathbb Z}}$

\[
\widetilde{\mathcal E}/T\simeq(\mathfrak S_3\times\widehat{\mathbb Z}^{\times})\underset{1/2}{\cdot}\Pi_\infty .
\]

D'autre part, on a

\[
\Gamma\mathrel{\rule{3em}{0.4pt}}\mathcal A
\tag*{(40)}
\]

hom dont l'image \inserted{$\Gamma'$} est d'indice fini (est-on $=$ surjectif) $\Gamma'$ est le groupe de Galois de la sous-extension de $\unsure{\overline{\mathbb Q}|\mathbb Q}/\unsure{\overline{\mathbb Q}|\mathbb Q}$ engendrée par toutes les racines $n$-ièmes de $\tfrac12$,

\margin{?}

% archive sheet number: 35
